\documentclass[11pt]{article}
\usepackage{amsmath,amssymb,booktabs,array,longtable,needspace}
\usepackage{hyperref}
\usepackage{komorebi-publication}
\ifdefined\pdfinfoomitdate\pdfinfoomitdate=1\fi
\ifdefined\pdftrailerid\pdftrailerid{}\fi
\ifdefined\pdfsuppressptexinfo\pdfsuppressptexinfo=15\fi
\renewcommand{\KomorebiTitle}{Arithmetic Packet Realization}
\renewcommand{\KomorebiSubtitle}{A worked visual atlas of constructions, obstructions, and the remaining conjecture}
\renewcommand{\KomorebiSeries}{Paper XI / Conjecture 66}
\renewcommand{\KomorebiKind}{Visual research companion}
\renewcommand{\KomorebiEdition}{Working atlas v0.5.0 / Publication Block 05 / 10 October 2026}
\renewcommand{\KomorebiShortTitle}{Paper XI / Conjecture 66 / Worked atlas}
\renewcommand{\KomorebiTitleSize}{32}
\renewcommand{\KomorebiTitleLeading}{37}
\KomorebiSetup
\newcommand{\F}{\mathbf F}
\newcommand{\C}{\mathbf C}
\newcommand{\Z}{\mathbf Z}
\newcommand{\StartPlate}[3]{\clearpage\phantomsection\label{plate:#1}\pdfbookmark[1]{#1 / #2}{#1}\noindent{\sffamily\small\bfseries\color{KomorebiMuted}WORKED PLATE #1}\par\vspace{5pt}{\sffamily\fontsize{21}{25}\selectfont\bfseries\color{KomorebiPine}#2\par}\vspace{6pt}{\sffamily\normalsize #3\par}}
\newcommand{\PlateSource}[3]{\par\noindent{\color{KomorebiPine}\rule{\linewidth}{.4pt}}\par{\small\sffamily\textbf{Evidence:} #1.\par\textbf{Companion source:} #2. PDF page #3 onward.}}
\begin{document}
\KomorebiCover
\pdfbookmark[0]{Reading the atlas}{reading}
\section*{A reading route through the argument}
This atlas accompanies \emph{Toward Arithmetic Packet Realization}, the 77-page cumulative research paper v0.4.0. Each plate combines one spacious mathematical graphic with an explicit claim boundary and a route back to that paper. The figure data are exact finite calculations or labeled dependency diagrams, not invented measurements.

\textbf{The central distinction.} A finite packet, an actual arithmetic host, and the same candidate's full local factors are different objects. The first two plates fix the quantifiers. The rest explain what can be constructed, what a particular obstruction rules out, and which stronger inputs are still required.

\begin{center}\renewcommand{\arraystretch}{1.6}
\begin{tabular}{@{}p{.15\linewidth}p{.76\linewidth}@{}}\toprule
Route & Mathematical purpose\\\midrule
\hyperref[plate:V01]{V01--V05} & The witness contract, balanced products, compact binders, and a complete finite corner obstruction.\\
\hyperref[plate:V06]{V06--V10} & Whole characters, trace depth, defining sequences, native quotient capacity, and enlarged-carrier recovery.\\
\hyperref[plate:V11]{V11--V19} & Conductor conventions, candidate filters, information loss, descent, Clifford dimensions, and occurrence/action obstructions.\\
\hyperref[plate:V20]{V20--V26} & Full finite extensions, section products, quadratic markings, the equal-capacity boundary, and transvection repair.\\
\hyperref[plate:V27]{V27--V32} & Genuine action correction, its even-order limit, honest phases, the missing native certificate, and the remaining endpoints.\\\bottomrule
\end{tabular}\end{center}

\textbf{Evidence labels.} An exact finite model verifies the displayed example. A conditional profile comparison retains the actual arithmetic construction hypotheses. A contract diagram records dependencies but does not supply their inputs. Ordinary general proofs remain in the companion paper. None of these labels denotes a new Lean run or external peer review.

The source line on every plate links to the cited physical PDF page when the atlas and research paper are kept in the same folder. The final concordance gives exact section and page routes. The original core reader trio is unchanged in this checkpoint.
\StartPlate{V01}{The three obligations}{Finite comparison, arithmetic existence, and characterization have different inputs.}
\begin{center}\includegraphics[width=\linewidth]{figures/V01.pdf}\end{center}
The source is a wild Weil atom $\rho$ with the supplied datum $d_\rho$. H1--H4 include the declared finite break decomposition, symplectic bridges, central extensions and phases, and coherent transports. They describe an admitted input. They do not assert that every possible source arrives with canonical data.

The middle box is an actual arithmetic witness $R(d_\rho,\pi)$. It carries the prescribed packet and recovers the same supercuspidal representation, including the central character and noncompact completion. The final arrow additionally needs H5 for that particular $\pi$ and the applicable characterization theorem.
\[
\forall\rho\in\mathcal A\quad\exists\pi\;R(d_\rho,\pi)
\qquad\hbox{is the general realization endpoint.}
\]
An arrow is a dependency, not a declaration that its input has been constructed. Restricted witnesses and finite comparison theorems remain useful even while this quantified endpoint is open. Historical obligation 66 is not a measure of the amount of mathematics still required.
\vfill\PlateSource{Contract diagram}{Part I, Section 1}{6}

\StartPlate{V02}{One candidate must satisfy both conditions}{Separate existential statements cannot be joined by changing the witness.}
\begin{center}\includegraphics[width=\linewidth]{figures/V02.pdf}\end{center}
Let $A=\{\pi:R(d,\pi)\}$ and $B=\{\pi:H5(\pi)\}$. The desired comparison uses a member of $A\cap B$. The diagram is an exact two-element countermodel: $A=\{\pi_1\}$ and $B=\{\pi_2\}$, with $\pi_1\ne\pi_2$. Each condition has a witness. Their conjunction has none.
\[
(\exists\pi\,R(d,\pi))\land(\exists\pi\,H5(\pi))
\quad\not\Rightarrow\quad
\exists\pi\,(R(d,\pi)\land H5(\pi)).
\]
There is a second pitfall. An implication about every realizer can be true when there are no realizers. A comparison interface whose structure already contains an existence field does not construct that field merely by using it.

These finite sets are logical illustrations, not a claim that the arithmetic realization and H5 classes are actually disjoint. Their purpose is to make a change of candidate visible before it enters a proof.
\vfill\PlateSource{Exact logical countermodel}{Part I, Section 1.2}{6}

\StartPlate{V03}{A balanced central product}{The shared center is identified once, while the orthogonal packet factors remain.}
\begin{center}\includegraphics[width=\linewidth]{figures/V03.pdf}\end{center}
Take two Heisenberg groups of order $3^3=27$. Their direct product has order $729$. Quotient by the central subgroup $\{(z,z^{-1})\}$ of order three. The balanced central product has order $3^5=243$, with a four-dimensional quotient over $\F_3$ and one common center.
\[
(H_3\times H_3)/\{(z,z^{-1})\}
\quad\text{has Heisenberg degree}\quad 3^2=9.
\]
The figure's order calculation is checked by enumerating every fibre of the multiplication model. Each fibre has exactly three elements. The representation degree agrees with the tensor product of two degree-three Heisenberg representations whose central characters match.

The general argument requires the actual factor extensions and commutation between distinct factors. A list of break labels does not supply either condition. Once those conditions hold, the balanced product removes the proposed extra mixed extension ambiguity. It does not construct an arithmetic host or identify independently fixed lifts.
\vfill\PlateSource{Exact finite model}{Part I, Section 3}{8}

\StartPlate{V04}{What a compact binder can do}{Compatible solutions in one fixed compact space can survive every precision.}
\begin{center}\includegraphics[width=\linewidth]{figures/V04.pdf}\end{center}
The points are all roots of $x^2-1$ modulo $3^n$, for $1\le n\le5$. The labels are canonical residue representatives. The arrows reduce a root modulo $3^n$ to a root modulo $3^{n-1}$. Both branches persist, corresponding to $1$ and $-1$ in $\Z_3$.

For the general binder, fix a closed $X\subseteq\mathcal O^m$ and continuous constraints $f_j$. The sets
\[
C_n=\{x\in X:f_j(x)\in\varpi^n\mathcal O\text{ for every }j\le n\}
\]
are closed and nested. If every $C_n$ is nonempty, compactness supplies a point satisfying every exact equation.

The example does not prove simultaneous finite solvability for the arithmetic problem. Even $x=0$ and $x=1$ have separately solvable congruences but no joint solution modulo three. Allowing the number of coefficients or pole bounds to grow also changes the compact model. Those missing hypotheses cannot be replaced by more isolated successful tests.
\vfill\PlateSource{Exact finite illustration}{Part I, Section 5; Part II, Section 19.1}{10}

\StartPlate{V05}{The complete cubic corner spectrum}{A finite obstruction can exclude a template without excluding every host.}
\begin{center}\includegraphics[width=\linewidth]{figures/V05.pdf}\end{center}
Rows run through the eight nonzero triples $(u,v,w)\in\F_3^3$. Columns run through all nine pairs $(p,q)$. A circle marks a solvable target in the specified corner. A cross marks an excluded target.
\[
u=cd,\quad v=ab,\quad w=bc,\quad
p=c(a-h),\quad q=b(d-h).
\]
Elimination gives the exact criterion $p+u=w(q+v)$. Thus precisely $24$ of the $72$ targets are admissible. The data enumerate all five coefficients over $\F_3$ and independently compare their image with this criterion.

The target $(1,1,1,2,2)$ becomes possible when the fifth coefficient $h$ is allowed. The target $(1,1,1,1,2)$ remains impossible. The inherited structural calculation explains why the larger corner has the same decisive leading condition. Its scope is that corner. The later arbitrary-exponent binder changes the construction rather than claiming that this negative result vanished.
\vfill\PlateSource{Exhaustive finite spectrum}{Part II, Section 19.2}{50}

\StartPlate{V06}{Do not truncate the full character}{Deep values, the full compact domain, and a noncompact generator are distinct certificates.}
\begin{center}\includegraphics[width=\linewidth]{figures/V06.pdf}\end{center}
A trace expression first describes a character on sufficiently deep congruence units. The full simple character also has values on the larger native unit domain. A full inducing representation retains further central and noncompact data. These are successive obligations, not three interchangeable notations for the same calculation.

The direct character construction transfers the whole preceding character, applies the actual permitted correction, and then checks the intersection needed for shallow extension. The continuation113 correction exhibits why noncompact normalization matters. In exponents of ninth roots of unity,
\[
\nu(\varpi)=\zeta_9,\qquad
\nu(\varpi_F/\varpi^3)=\zeta_3^2
\quad\Longrightarrow\quad \nu(\varpi_F)=1.
\]
Setting $\nu(\varpi)=1$ instead gives $\zeta_3^2$, changing the central character. The accompanying finite check verifies this root-of-unity identity only. The full unit character, intertwining, and type-extension statements retain the arithmetic hypotheses developed in the progress paper.
\vfill\PlateSource{Domain diagram}{Part I, Sections 7.5--8.5; Part II, Section 18}{15}

\StartPlate{V07}{Depth can make native trace one invisible}{The valuation inequality is the mechanism, not an assertion that the trace is zero.}
\begin{center}\includegraphics[width=\linewidth]{figures/V07.pdf}\end{center}
In $\F_2((t))$, let the additive character read the coefficient of $t^0$. Model the trace coefficient by $t^\tau$, and let the variable range over $t^r\mathcal O$. A circle means every such product has positive valuation and is invisible to the character:
\[
\tau+r\ge1.
\]
At a cross, $t^{-\tau}$ belongs to the allowed domain and produces the visible product $1$. This gives an explicit negative control, rather than treating failure of a sufficient bound as proof of visibility without a witness.

The native trace-one row is $\tau=0$. It is invisible at every positive unit depth in the displayed range. Native trace one, zero trace after amplification, and the trace of the correction remain different quantities in the arithmetic recursion. The exact floor in the correction depth is retained in the written proof. This monomial picture explains the inequality but does not replace the full shallow-character extension argument.
\vfill\PlateSource{Exact Laurent-monomial model}{Part I, Section 8}{18}

\StartPlate{V08}{Two defining-sequence examples}{The recursion records levels and field refinements, not just its final matrix size.}
\begin{center}\includegraphics[width=\linewidth]{figures/V08.pdf}\end{center}
The odd example uses $p=3$, $d=2$, and $(h_1,h_2)=(1,2)$. Its displayed levels are $(24,4)$. The dyadic example uses $d=3$ and $(h_1,h_2,h_3)=(1,3,5)$, giving $(60,18,3)$. Each position refers to that example's own defining sequence.
\[
\nu_j=p^{d-j}(p+1)h_{d-j+1}
\quad\text{or}
\quad\nu_j=2^{d-j}3h_{d-j+1}.
\]
The two plotted lines are not empirical fits or samples of one common arithmetic family. They make the inherited level formulas legible. In the odd example, $24>3\cdot4>2\cdot4$ illustrates the separated scales used in refinement.

Supplying the levels does not prove that the literal corrected matrix generates the asserted maximal field. That conclusion uses the filtered corestriction lift and the precise refinement theorem. In the dyadic argument, separability from native trace one is established only after the full field degree has been obtained.
\vfill\PlateSource{Inherited formulas, exact evaluation}{Part I, Sections 7.1--7.4 and 8.1--8.3}{12}

\StartPlate{V09}{Where the original dyadic quotient runs out of room}{The rank-eight case is the first equality, not a low-rank surplus.}
\begin{center}\includegraphics[width=\linewidth]{figures/V09.pdf}\end{center}
For $d$ source planes, the binary source quotient has dimension $2d$. In the full dyadic profile the specified original maximal-type quotient has dimension $D=2^d-2$. The first five values are shown without normalizing away their difference.
\[
\begin{array}{c|ccccc}
d&1&2&3&4&5\\
2d&2&4&6&8&10\\
2^d-2&0&2&6&14&30
\end{array}
\]
Ranks two and four are too small in that original quotient. Rank eight has exactly six dimensions for six required dimensions. Only from rank sixteen is there a strict surplus in this family.

The later low-rank repair uses a larger actual assembly carrier and recovers the same representation. It does not change the first two numbers in this table. At equality, a fixed native quadratic refinement can still obstruct attachment. With spare dimension, the free quadratic embedding problem is easier, but prescribed actions, lifts, and transports remain separate constraints.
\vfill\PlateSource{Exact capacity formula}{Part I, Section 9; Part III, Section 24.2}{21}

\StartPlate{V10}{A larger inducing carrier can return the same object}{Induction in stages changes the intermediate group without changing the result.}
\begin{center}\includegraphics[width=\linewidth]{figures/V10.pdf}\end{center}
The diagram is a finite analogue, not the actual local-field construction. Use $H=C_2\subset M=C_6\subset G=C_{12}$ and a character of $H$. Induction to $M$ has dimension three, and induction onward to $G$ has dimension six. Direct induction has the same six cosets, represented by $i+2j$ with $i\in\{0,1\}$ and $j\in\{0,1,2\}$.
\[
\operatorname{Ind}_M^G\operatorname{Ind}_H^M\lambda
\simeq\operatorname{Ind}_H^G\lambda.
\]
The arithmetic recovery uses the corresponding compact-induction transitivity with the original full inducing representation $\Lambda$. Its larger intermediate carrier must be an actual permitted subgroup, and the original central character and noncompact extension must be retained.

The rank-two and rank-four proofs supply additional group, normalization, and character-orbit arguments before transitivity applies. The finite coset calculation isolates the recovery mechanism. It neither constructs those arithmetic groups nor proves irreducibility of a new local-field representation.
\vfill\PlateSource{Exact finite analogue of recovery}{Part I, Section 10}{22}

\StartPlate{V11}{Three conductor columns are not interchangeable}{The scalar summand changes an epsilon exponent even when the ordinary conductor is unchanged.}
\begin{center}\includegraphics[width=\linewidth]{figures/V11.pdf}\end{center}
The ordinary conductor is $a(V)=\dim V-\dim V^{I_F}+\operatorname{Sw}(V)$. For the additive character $\psi_1$ trivial on $\mathfrak p_F$ but not on $\mathcal O_F$, the stated normalization gives $f_{\psi_1}(V)=a(V)-\dim V$.
\[
\begin{array}{c|rrr}
V&\operatorname{Sw}(V)&a(V)&f_{\psi_1}(V)\\
\rho&24&33&24\\
\operatorname{End}\rho&152&232&151\\
\operatorname{Ad}\rho&152&232&152
\end{array}
\]
Here $\operatorname{Ad}\rho$ means the trace-zero summand. The scalar summand in $\operatorname{End}\rho=\C\oplus\operatorname{Ad}\rho$ has zero ordinary conductor, but it contributes one to the dimension subtracted in the epsilon normalization.

The plot keeps the three definitions separate even when values happen to coincide. An exponent is still not the full epsilon factor. The arithmetic conductor inputs and additive-character convention are inherited from the audited progress paper, not freshly established by drawing these numbers.
\vfill\PlateSource{Exact normalization worksheet}{Part III, Section 22.1}{58}

\StartPlate{V12}{The principal-period filter}{The period belongs to the actual order, and the level must be integral.}
\begin{center}\includegraphics[width=\linewidth]{figures/V12.pdf}\end{center}
Under the stated actual-type conductor formula, a rank-$n$ source with Swan conductor $w$ requires $nr=ew$, where the principal period satisfies $e\mid n$. Therefore
\[
\frac{n}{\gcd(n,w)}\mid e\mid n,
\qquad r=\frac{ew}{n}.
\]
For $(n,w)=(9,24)$, period one would demand the nonintegral level $8/3$. Periods three and nine leave $(e,r)=(3,8)$ and $(9,24)$. The cross is excluded numerically. The circles pass only this integrality condition.

A period-three candidate can still fail the independent torsion-number test when its actual maximal parameter field has residue degree three. The prescribed wild-irreducible source has no nontrivial unramified self twist. Neither a convenient lattice period nor a lower critical level can be substituted for the actual period and leading level in the formula. A necessary numerical pass supplies no simple stratum or character occurrence by itself.
\vfill\PlateSource{Exact necessary filter}{Part III, Section 22.2--22.3}{58}

\StartPlate{V13}{The rank-nine source has a visible break count}{The self-pair conductor comes from the unchanged source and its actual independence assumptions.}
\begin{center}\includegraphics[width=\linewidth]{figures/V13.pdf}\end{center}
For the recorded source $\rho=A_1\otimes A_2$, each degree-three Heisenberg factor contributes one trivial conjugation character and eight nontrivial characters, at break one or two respectively. Under the original independence hypotheses, the tensor product contributes one trivial character, eight at break one, and seventy-two at break two.
\[
\operatorname{Sw}(\operatorname{End}\rho)=8+2\cdot72=152,
\qquad a(\operatorname{End}\rho)=81-1+152=232.
\]
The last subtraction uses wild irreducibility: the invariant endomorphisms are exactly the scalars. It does not require Frobenius to fix every individual character after the splitting extension.

The histogram is computed by enumerating the $9\times9$ pairs and taking the larger of the two unequal breaks. The arithmetic assertions that these are the source's genuine characters, that their breaks are as stated, and that the quotient factors are independent remain the inherited source hypotheses. Their role is exposed rather than replaced with a finite counting experiment.
\vfill\PlateSource{Exact character-count model}{Part III, Section 23.1}{59}

\StartPlate{V14}{A conductor filter is not an identification}{The critical-four candidate passes one invariant; it is not thereby identified with the source.}
\begin{center}\includegraphics[width=\linewidth]{figures/V14.pdf}\end{center}
The two points use the recorded period-nine actual-order certificate. It gives leading level $24$, an even critical level $c$, index exponent $I=c+64$, and quotient dimension $D=8$. With the stipulated measures and torsion number one, the self-pair conductor is
\[
a(\pi^\vee\times\pi)=2c+224.
\]
At $c=10$ this is $244$, contradicting the unchanged source value $232$. At $c=4$ it is $232$, so this necessary comparison no longer excludes the candidate. The figure shows only these two conditional evaluations, not an asserted family of actual constructions for all intermediate values of $c$.

The order equations are substantive inputs. The pair $(24,c)$ does not imply them. Passing the horizontal line neither constructs the corrected stratum nor establishes occurrence of the complete prescribed character. Full pair-factor equality for the same candidate remains a stronger obligation than this conductor comparison.
\vfill\PlateSource{Conditional profile comparison}{Part III, Section 23.2--23.3}{60}

\StartPlate{V15}{Equal weighted sums can conceal different capacity}{The conductor formula has already discarded the parity information needed for the quotient.}
\begin{center}\includegraphics[width=\linewidth]{figures/V15.pdf}\end{center}
The full rank-eight profile has centralizer dimensions $(8,4,2,1)$ and levels $(60,18,3)$. The skipped profile has dimensions $(8,4,1)$ and levels $(60,13)$. The bar segments are the products $\Delta_j\nu_j$. Both sum to $279$, giving ordinary self-pair conductor $63+279=342$.
\[
(I,D)_{\rm full}=(133,6),\qquad (I,D)_{\rm skipped}=(134,4).
\]
The compensation $2\cdot133+6=2\cdot134+4$ explains the equality. The identity $2\lceil\nu/2\rceil+\mathbf1_{2\mid\nu}=\nu+1$ removes the parity distribution from the weighted conductor sum, while $D$ retains it.

Under their inherited arithmetic construction hypotheses, both profiles also have standard conductor $68$ and torsion number one. Yet only the full profile has room for the six-dimensional source quotient in this specified host. The new tests verify the grade consequences. They do not freshly construct the skipped-profile parameter or turn three agreeing integers into H5.
\vfill\PlateSource{Exact information-loss example}{Part III, Section 24}{62}

\StartPlate{V16}{Central normalization of a descent fibre}{A power map can be bijective in one fibre and collapse an entire family in another.}
\begin{center}\includegraphics[width=\linewidth]{figures/V16.pdf}\end{center}
A descent fibre that is a torsor for a cyclic character group $X$ of order $m$ is normalized by the central-character power map $\alpha\mapsto\alpha^n$. Additive exponents identify it with $x\mapsto nx\pmod m$.

On the left, $n=9$ and $m=4$. The map is the identity modulo four, so every attainable central-character ratio selects one twist. On the right, $n=3$ and $m=3$. Every exponent maps to zero, so the fixed-center condition leaves all three twists available.
\[
\gcd(n,m)=1\quad\Longrightarrow\quad
\text{unique member for each attainable central ratio.}
\]
The two examples correspond to different source branches. The rank-nine positive-characteristic descent and the rank-three continuation113 fixed-center ambiguity are not one candidate under two names. Uniqueness in the coprime fibre does not prove a native character occurs in that candidate. Ambiguity in the cubic fibre does not disappear because a scalar obstruction has been removed.
\vfill\PlateSource{Exact finite character-group model}{Part II, Sections 17.3 and 18.4; Part III, Section 25.3}{46}

\StartPlate{V17}{Clifford dimensions keep the radical and orbit factors}{An abstract packet count must match an actual restriction problem.}
\begin{center}\includegraphics[width=\linewidth]{figures/V17.pdf}\end{center}
For a finite abelian projective quotient $A$, let $R$ be the radical of its commutator pairing. An irreducible projective multiplicity has degree
\[
(\dim M)^2=|A/R|.
\]
The first example uses $A=\F_2^3$, with the standard symplectic pairing on the first two coordinates and the third coordinate radical. Thus $|A|=8$, $|R|=2$, and $\dim M=2$. Suppressing the radical would incorrectly produce $\sqrt8$.

The second example retains the inertia orbit. The two-dimensional irreducible of $S_3$ restricts to two distinct nontrivial characters of $C_3$. Each multiplicity is one, but the orbit has size two. Along a genuine compatible chain, every orbit factor and every projective multiplicity must be retained.

Independent commuting projective factors on the same carrier multiply their degrees. Representing one Weyl plane twice on the same space does not create independent commuting factors. The finite matrix checks include the $S_3$ relations, while the general restriction argument remains in the paper.
\vfill\PlateSource{Exact projective and orbit models}{Part II, Section 15}{41}

\StartPlate{V18}{The noncompact scalar obstruction survives compact agreement}{Changing the full character removes one obstruction, not the occurrence problem.}
\begin{center}\includegraphics[width=\linewidth]{figures/V18.pdf}\end{center}
In continuation113, the original $\theta_0$ is compatible with the tested compact action but produces the wrong scalar at a noncompact generator. The retained intertwining and maximal-type arguments yield $\operatorname{Hom}_{H^1}(\theta_0,\pi_F^0)=0$ for the prescribed source-defined candidate.

The correction powers $\nu^a$, for $a=0,1,2$, leave generator scalars $\zeta_3,1,\zeta_3^2$. Only the middle correction passes this test. It changes the full character, rather than merely changing a basis. Its uniformizer value must also preserve the central normalization shown on V06.
\[
\operatorname{Hom}_{H^1}(\theta',\pi_F^0)\ne0
\quad\text{is still not established.}
\]
The new atlas recomputes the exact scalar exponents. It does not rerun the earlier local-field matrix programs or supply a new occurrence proof. Choosing some other representation that contains $\theta'$ would abandon the fixed-candidate question. Rescaling an intertwiner cannot remove a nontrivial scalar between two representations whose common compact restriction is irreducible.
\vfill\PlateSource{Inherited arithmetic obstruction with exact scalar model}{Part II, Section 18}{47}

\StartPlate{V19}{Abstract actions and native normalizers}{Abstract lifting and arithmetic realization are different existence statements.}
\begin{center}\includegraphics[width=\linewidth]{figures/V19.pdf}\end{center}
The linear map $A(x,y)=(-y,x)$ on $\F_3^2$ is symplectic and has order four. Thus there is a genuine nontrivial abstract $C_4$ action. In contrast, every homomorphism $C_4\to C_{3^r}$ is trivial: an image $x$ satisfies $4x=0$, and multiplication by four is invertible modulo $3^r$.

The arithmetic obstruction uses its stated native-normalizer quotient and inner-action requirements. When that permitted quotient is a $p$-group, a prime-to-$p$ subgroup cannot implement the desired nontrivial action through it. The cyclic groups in the drawing isolate the elementary coprimality mechanism. They are not asserted to be the entire normalizer of every type.

Finite central-extension actions can still lift under their own coprimality hypotheses, and projective intertwiners can sometimes be linearized. Neither result supplies the missing arithmetic conjugating matrices. An obstruction to this normalizer mechanism therefore remains compatible with success through a different allowed host or a different stated attachment contract.
\vfill\PlateSource{Exact coprimality analogue}{Part I, Sections 4 and 11}{9}

\StartPlate{V20}{The finite subgroup remembers its squares}{Dihedral and quaternion markings survive inside the same complex scalar extension.}
\begin{center}\includegraphics[width=\linewidth]{figures/V20.pdf}\end{center}
Let $X^2=Z^2=I$ and $XZ=-ZX$. The subgroup generated by $X,Z$ is dihedral of order eight. The subgroup generated by $iX,iZ$ is quaternion of order eight. Both have the same quotient polarization and commutator $-I$.
\[
q_D(x,y)=xy,\qquad q_Q(x,y)=xy+x+y.
\]
Their lift squares differ. The order columns count elements of exact order two and four, excluding the identity. They distinguish the finite groups even though adjoining all complex scalars gives the same ambient scalar extension.

For a fixed two-element center, a center-fixing extension map over a prescribed symplectic map must preserve the quadratic refinement. With complex scalar center, square ratios have square roots, but this freedom does not erase a separately prescribed finite subgroup. Transporting a quaternion marking into the common scalar extension is valid. Identifying it with the independently fixed dihedral marking is not.
\vfill\PlateSource{Exact finite-group enumeration}{Part III, Section 26.1--26.2}{67}

\StartPlate{V21}{A square table does not specify an entire section}{Changing one section value can preserve every square and still change products.}
\begin{center}\includegraphics[width=\linewidth]{figures/V21.pdf}\end{center}
Start from the ordered section $\sigma(x,y)=X^xZ^y$ in the dihedral example. Change only $\sigma(e+f)$ by a central sign. Write the sign function as $a(v)=\mathbf1_{v=e+f}$. The displayed entries are the resulting cocycle differences
\[
\Delta c(v,w)=a(v)+a(w)+a(v+w)\in\F_2.
\]
Every diagonal entry is zero, so every square is unchanged. But $\Delta c(e,f)=1$. A map fixing the specified generating lifts cannot also implement the changed assignment at $e+f$. The cross marks a nonzero product discrepancy, not failure of the two extensions to be isomorphic in some other way.

This is why the native certificate needs all specified section products, not only a quadratic table. When the entire section is transported through one homomorphism, those products are retained automatically. When the target section is independently prescribed, its pointwise assignments must be checked. The distinction concerns what is fixed, not a contradiction in extension theory.
\vfill\PlateSource{Exact cocycle counterexample}{Part III, Section 26.3}{68}

\StartPlate{V22}{Two different marking contracts}{A positive attachment and a native obstruction can both be correct.}
\begin{center}\includegraphics[width=\linewidth]{figures/V22.pdf}\end{center}
In the transported-marking contract, the source marking is supplied. After a permitted symplectic embedding and scalar-extension isomorphism have been chosen, the host marking is their image. A homomorphism fixes the represented center and carries squares, products, and specified lifts together.

In the independently fixed native contract, the target finite subgroup and its refinement are already given. The comparison must hit that subgroup with its existing marking. A scalar rescaling that changes it may be disallowed.
\[
\text{choose a compatible target marking}
\quad\ne\quad
\text{match this prescribed target marking}.
\]
The restricted one-vertex, trivial-stabilizer attachment retains its actual-type, capacity, full-character, and recovery hypotheses. It should not acquire an extra native-marking premise merely because a stronger problem is also studied. Conversely, that positive theorem cannot certify the stronger problem without computing the additional datum. The two branches in the figure remain separate throughout the atlas.
\vfill\PlateSource{Contract comparison}{Part III, Sections 26.4 and 27.3}{68}

\StartPlate{V23}{Two six-dimensional refinements with different signs}{The alternating form and dimension do not determine the Arf invariant.}
\begin{center}\includegraphics[width=\linewidth]{figures/V23.pdf}\end{center}
Each grid contains all $64$ binary vectors in dimension six. Columns encode the lower three bits and rows the upper three bits, with row zero at the top. Circles have quadratic value zero. Crosses have value one. The two forms share the standard polarization:
\[
q_+=x_1y_1+x_2y_2+x_3y_3,
\qquad q_-=q_++x_1+y_1.
\]
Their Gauss sums are $+8$ and $-8$. Consequently the zero counts, including the zero vector, are $(64+8)/2=36$ and $(64-8)/2=28$. The diagrams use exact binary evaluation at every vector.

These are explicit abstract refinements, not a claim that either has been computed on the actual rank-eight arithmetic candidate. The example shows the datum still required when the host marking is independently fixed. At equal capacity, an Arf mismatch has no orthogonal complement in which to be absorbed.
\vfill\PlateSource{Exhaustive 64-vector displays}{Part III, Section 27.1}{69}

\StartPlate{V24}{The equal-capacity test is sharp}{One spare quadratic plane removes the free Arf obstruction, not every attachment constraint.}
\begin{center}\includegraphics[width=\linewidth]{figures/V24.pdf}\end{center}
Fix a nondegenerate source of dimension six. The table varies the host dimension and its supplied Arf invariant. The embedding may be freely chosen. A four-dimensional host cannot contain the source. At dimension six, the map is an isometry and the Arf signs must agree. At dimension eight, the complementary plane can have either sign.
\[
\operatorname{Arf}(U)
=\operatorname{Arf}(W)+\operatorname{Arf}(V)
\quad\text{for }W\simeq V\perp U.
\]
This is the finite classification argument, not a numerical search over arithmetic matrices. It combines with the original full dyadic capacity formula to isolate the rank-eight equality case.

The word ``free'' matters. A prescribed depth placement, an already fixed quotient embedding, or a stabilizer-equivariance condition may forbid the required change of subspace. The table assumes the actual quadratic host has been supplied. It does not compute that host's native marking, choose its lifts, or authorize modifications to the transport diagram.
\vfill\PlateSource{Finite classification criterion}{Part III, Section 27.2--27.3}{69}

\StartPlate{V25}{A single transvection repairs an isotropic defect}{The repair is explicit, and the previously correct line stays fixed.}
\begin{center}\includegraphics[width=\linewidth]{figures/V25.pdf}\end{center}
Take $q(x,y)=xy$ and $q_0(x,y)=xy+y$. Their discrepancy is $B(e,-)$, with $e=(1,0)$ and $q(e)=0$. The transvection is
\[
T_e(x,y)=(x+y,y),\qquad q_0(T_e v)=q(v).
\]
The arrows show its action on every vector. Zero and $e$ remain fixed. The vectors $f$ and $e+f$ exchange places. Each arrow therefore lands on the same quadratic value shown at its source.

The original identity map did not work: $q(f)=0$ while $q_0(f)=1$. Matching Arf invariants guarantees a correction here, not the validity of an unchanged map. More generally, a defect vanishing on a subspace $U$ gives a defect vector in $U^\perp$, so the transvection fixes $U$ pointwise.

The correction fixes quotient vectors, not automatically independently fixed lifts above them. Its use in an arithmetic diagram requires permission to change the embedding and a separate check of the relevant section and phase data.
\vfill\PlateSource{Exact four-vector correction}{Part III, Section 28.1--28.3}{71}

\StartPlate{V26}{The defect vector supplies a decision procedure}{The same formula displays both the successful correction and its genuine obstruction.}
\begin{center}\includegraphics[width=\linewidth]{figures/V26.pdf}\end{center}
For the fixed hyperbolic plane $q=xy$, every refinement with the same polar form is $q_0=q+B(a,-)$ for one vector $a$. The table lists all four defect vectors. Its mismatch column counts the vectors at which $q_0(T_av)\ne q(v)$.
\[
q_0(T_av)=q(v)+B(a,v)q(a),\qquad
\operatorname{Arf}(q_0)+\operatorname{Arf}(q)=q(a).
\]
When $q(a)=0$, the correction works everywhere. For $a=e+f$, the Arf sign changes and two vectors still fail. This is not an unlucky choice of basis. In equal dimension, the two forms are not isometric.

In a larger host this test applies to the pullback along the supplied embedding. A different permitted subspace may behave differently. Under genuine supplied actions preserving both refinements and an initially equivariant embedding, the defect vector is fixed by the group. The same transvection is then equivariant. Those initial actions are hypotheses, not outputs of this linear solve.
\vfill\PlateSource{Exact defect enumeration}{Part III, Section 28}{71}

\StartPlate{V27}{An odd-order action discrepancy can be summed away}{The correction acts on a supplied finite extension, not on an unconstructed arithmetic host.}
\begin{center}\includegraphics[width=\linewidth]{figures/V27.pdf}\end{center}
Let $C_3$ act on $A=\F_2^2$ by $g(x,y)=(y,x+y)$. Use the cocycle values $\delta_1=0$, $\delta_g=e+f$, and $\delta_{g^2}=f$. Their sum is $u=e$. The final column displays the corrected discrepancy, which is zero at every group element.
\[
\delta_{gh}=\delta_g+g\delta_h,
\qquad u=\sum_h\delta_h,
\qquad \delta_g=u+gu.
\]
The last identity follows by summing the cocycle equation and using that the group order is odd in characteristic two. Replacing the comparison by the vertical sign correction $u$ makes it equivariant.

In the extension theorem, the source and target already have genuine actions and an initial center-fixing extension isomorphism. The cocycle lives in the character space $\operatorname{Hom}(W,\F_2)$. The diagram is an exact module example of that calculation, not a new construction of those actions. A vertical correction may also be forbidden when particular lifts must remain unchanged.
\vfill\PlateSource{Exact C3-module model}{Part III, Section 29.1--29.2}{73}

\StartPlate{V28}{The even-order negative control}{A nonzero cocycle survives when every permitted coboundary vanishes.}
\begin{center}\includegraphics[width=\linewidth]{figures/V28.pdf}\end{center}
Let $C_2=\{1,s\}$ act trivially on the nonzero module $A=\F_2^2$. The assignment $\delta_1=0$, $\delta_s=e$ satisfies the cocycle equation, since $e+e=0$. For every proposed correction $u$,
\[
u+su=u+u=0,
\qquad u+\delta_s+su=e.
\]
The four rows exhaust the module. None removes the discrepancy. Thus the odd-order hypothesis of the preceding plate is mathematically active.

This does not say that every even-order action problem is obstructed. It supplies one exact counterexample to an unrestricted averaging claim. Other actions and cocycles may admit corrections. In the arithmetic problem, it remains necessary to compute the actual discrepancy and respect the permitted transformations. Declaring every projective or equivariance mismatch removable would erase precisely the information this example retains.
\vfill\PlateSource{Exact obstruction model}{Part III, Section 29.2}{74}

\StartPlate{V29}{Honest phases remain a separate layer}{An action on the extension group is not yet the specified action on its representation.}
\begin{center}\includegraphics[width=\linewidth]{figures/V29.pdf}\end{center}
Let $g,h\in\{0,1,2\}$ represent $C_3$, and take the one-dimensional operators $U_g=\zeta_9^g$. Their multiplication has scalar discrepancy
\[
U_gU_h=\zeta_9^{c(g,h)}U_{g+h\bmod3},
\qquad c(g,h)=3\left\lfloor\frac{g+h}{3}\right\rfloor.
\]
The figure records its exponents modulo nine. Replacing $U_g$ by $\zeta_9^{-g}U_g$ gives an honest action. Yet the three honest characters $g\mapsto\zeta_3^{kg}$ are still distinct choices.

This scalar toy model isolates two separate questions: whether a projective multiplier can be removed, and which honest linearization the contract prescribes. It is not the phase system of the actual arithmetic candidate. The finite extension-action correction on V27 addresses a different discrepancy, in the group automorphisms. Canonical extension phases, Gauss conventions, and coherent transports must still be checked on the same full representation.
\vfill\PlateSource{Exact scalar projective toy model}{Part I, Section 4.2; Part III, Section 29.3}{9}

\StartPlate{V30}{The missing native certificate has explicit fields}{Question marks are not zeroes and must not be filled with a convenient abstract form.}
\begin{center}\includegraphics[width=\linewidth]{figures/V30.pdf}\end{center}
When the stronger contract independently fixes a native rank-eight marking, the certificate must name the actual candidate, the relevant finite subgroup and represented center, and the six genuine lifts of a symplectic basis. Their squares determine six bits only after those lifts and the two-element central representation have been established.
\[
q_W(v_i)\in\F_2\quad\text{must come from the actual lift square.}
\]
All prescribed section products must also be verified. Genuine extension actions do not by themselves choose the required honest representation phases. The certificate must state which maps or lifts may change and verify the resulting transport conditions.

No native square bit is supplied by this atlas. The data file records these fields as null, and the tests ensure they remain uncomputed rather than silently assigning zero. This is a conditional stronger obligation. It is not an extra premise inserted into the already valid transported-marking contract, where the finite marking is carried through the chosen comparison map.
\vfill\PlateSource{Unresolved arithmetic certificate}{Part III, Section 30.2}{75}

\StartPlate{V31}{Keep the two construction routes distinct}{A source-defined candidate and a type-defined candidate require an identification theorem.}
\begin{center}\includegraphics[width=\linewidth]{figures/V31.pdf}\end{center}
The monomial route begins with a full inducing character and an admissible tower. Under the stated classical induction and descent inputs it produces a source-defined candidate. Solvability of an ambient finite group does not make every proposed subgroup admissible for a prime-cyclic tower.

The direct route begins with an explicit parameter and type construction. Its refinement and full-character arguments produce an actual candidate under a restricted contract. Even when both routes match rank, standard conductor, torsion number, and self-pair conductor, those invariants do not identify their outputs.
\[
\pi_{\rm AI}\simeq\pi_{\rm type}
\quad\text{needs its own argument on the same source and field.}
\]
The positive-characteristic rank-nine source-defined route and the characteristic-zero continuation113 example remain separate cases. Neither is silently substituted for the other. The complete Henniart--Lemaire text was not freshly accessible in Block 4, so the detailed inherited theorem applications keep that audit boundary. This visual synthesis adds no new classical-source certification.
\vfill\PlateSource{Source-lineage diagram}{Part II, Sections 14 and 17; Part III, Section 25}{39}

\StartPlate{V32}{The current frontier}{Completed finite and restricted work remains reusable without relabeling the open endpoints.}
\begin{center}\includegraphics[width=\linewidth]{figures/V32.pdf}\end{center}
The progress record now has a distinct visual reading route through the restricted constructors, finite comparison mechanisms, exact obstructions, and conditional corrections. These are substantive results and reusable arguments. Their existence does not make the remaining arithmetic statements interchangeable.

General realization requires an actual common host with the admitted actions and transports. Same-candidate H5 requires the characterizing factors for that very realizer. The corrected-character problem requires occurrence in the original specified induction candidate. An independently fixed native marking, when required, needs the actual certificate on V30.

This atlas is Publication Block 5. It leaves the 77-page research paper byte-preserved at v0.4.0. The original 109-page monograph, 51-page visual guide, and 234-page Lean guide are inherited baselines awaiting coordinated revision. No fresh Lean compilation or external peer review is claimed. The versioned ledger records these separate statuses and the exact next editorial step rather than assigning a conjecture-completion percentage.
\vfill\PlateSource{Versioned claim map}{Part I, Section 12; Part II, Section 20; Part III, Section 30}{29}

\clearpage\pdfbookmark[0]{Plate concordance}{concordance}
\section*{Plate concordance}
Page numbers below are physical pages of the cumulative v0.4.0 PDF, not its reused printed page labels. Each plate also has a direct companion link.
\begin{center}\small\renewcommand{\arraystretch}{1.22}
\begin{tabular}{@{}p{.08\linewidth}p{.67\linewidth}r@{}}\toprule
Plate & Argument & PDF page\\\midrule
\hyperref[plate:V01]{V01} & The three obligations & 6 \\
\hyperref[plate:V02]{V02} & One candidate must satisfy both conditions & 6 \\
\hyperref[plate:V03]{V03} & A balanced central product & 8 \\
\hyperref[plate:V04]{V04} & What a compact binder can do & 10 \\
\hyperref[plate:V05]{V05} & The complete cubic corner spectrum & 50 \\
\hyperref[plate:V06]{V06} & Do not truncate the full character & 15 \\
\hyperref[plate:V07]{V07} & Depth can make native trace one invisible & 18 \\
\hyperref[plate:V08]{V08} & Two defining-sequence examples & 12 \\
\hyperref[plate:V09]{V09} & Where the original dyadic quotient runs out of room & 21 \\
\hyperref[plate:V10]{V10} & A larger inducing carrier can return the same object & 22 \\
\hyperref[plate:V11]{V11} & Three conductor columns are not interchangeable & 58 \\
\hyperref[plate:V12]{V12} & The principal-period filter & 58 \\
\hyperref[plate:V13]{V13} & The rank-nine source has a visible break count & 59 \\
\hyperref[plate:V14]{V14} & A conductor filter is not an identification & 60 \\
\hyperref[plate:V15]{V15} & Equal weighted sums can conceal different capacity & 62 \\
\hyperref[plate:V16]{V16} & Central normalization of a descent fibre & 46 \\
\bottomrule\end{tabular}\end{center}
\subsection*{Scope of the source routes}
The routes identify where the argument is developed in the supplied research paper. They do not reclassify an inherited source audit as a fresh one. The original arithmetic notes, negative controls, and classical-input qualifications remain in the preserved checkpoint.

The new plate files, exact model data, Python figure generator, and test runner are editable. The 77-page research paper is retained unchanged. Its Block 2 portion remains a preserved reader because the complete earlier editable archive has not been recovered.

\clearpage\pdfbookmark[0]{Concordance and replay boundary}{replay}
\section*{Concordance and replay boundary}
\begin{center}\small\renewcommand{\arraystretch}{1.22}
\begin{tabular}{@{}p{.08\linewidth}p{.67\linewidth}r@{}}\toprule
Plate & Argument & PDF page\\\midrule
\hyperref[plate:V17]{V17} & Clifford dimensions keep the radical and orbit factors & 41 \\
\hyperref[plate:V18]{V18} & The noncompact scalar obstruction survives compact agreement & 47 \\
\hyperref[plate:V19]{V19} & Abstract actions and native normalizers & 9 \\
\hyperref[plate:V20]{V20} & The finite subgroup remembers its squares & 67 \\
\hyperref[plate:V21]{V21} & A square table does not specify an entire section & 68 \\
\hyperref[plate:V22]{V22} & Two different marking contracts & 68 \\
\hyperref[plate:V23]{V23} & Two six-dimensional refinements with different signs & 69 \\
\hyperref[plate:V24]{V24} & The equal-capacity test is sharp & 69 \\
\hyperref[plate:V25]{V25} & A single transvection repairs an isotropic defect & 71 \\
\hyperref[plate:V26]{V26} & The defect vector supplies a decision procedure & 71 \\
\hyperref[plate:V27]{V27} & An odd-order action discrepancy can be summed away & 73 \\
\hyperref[plate:V28]{V28} & The even-order negative control & 74 \\
\hyperref[plate:V29]{V29} & Honest phases remain a separate layer & 9 \\
\hyperref[plate:V30]{V30} & The missing native certificate has explicit fields & 75 \\
\hyperref[plate:V31]{V31} & Keep the two construction routes distinct & 39 \\
\hyperref[plate:V32]{V32} & The current frontier & 29 \\
\bottomrule\end{tabular}\end{center}
\subsection*{What is freshly executable}
The atlas suite contains 28 test methods. It checks the plotted finite data and deliberate negative controls, including the 72-target corner, conductor translations, cocycle products, all 64 vectors in each six-dimensional display, and explicit correction formulas. The delivered replay also runs the 23-method Block 4, 15-method Block 3, and 12-method Block 1 suites. These scopes are reported separately rather than added into a measure of theorem completion.

No direct local-field matrix constructor, historical extra arithmetic program, Lean compiler, or external proof checker is run by that command. The executable native-certificate fields remain explicitly uncomputed. Full build and replay instructions are in \texttt{BUILD.md}. The versioned ledger and \texttt{CONTINUE\_HERE.md} specify the next checkpoint.
\end{document}
